% ---------------------------------------------------------------------------
% DRAFT SEGMENT — localization test stimuli
% Standalone; not \input anywhere yet. Stitch into Methods later.
% Backing analysis + parameter sweeps: docs/stimulus_spectral_variation.md
% Code: experiment/localization/localization_helpers/stimulus.py
% ---------------------------------------------------------------------------

\subsection{Localization stimuli}

Localization stimuli were 225\,ms trains of five 25\,ms pink-noise bursts
separated by 25\,ms silent gaps (10\,ms raised-cosine onset and offset ramps,
5\,ms at the gap edges), cut from a single continuous noise realization. In the
fixed-spectrum condition the train was used unmodified, so its expected
spectrum was the same on every trial and only the realization noise varied
(SD of the 1/6-octave-smoothed log spectrum: 0.4\,dB across trains, 0.8\,dB
across the bursts within a train).

In the variable-spectrum condition each train was recolored, before gapping, by
a random spectral envelope drawn anew on every trial and shared by all five
bursts. Natural sound sources differ from one another mainly in broad spectral
coloration, which the auditory system must discount in order to read the finer,
elevation-dependent pattern imposed by the pinna; the envelope was built to
impose exactly that kind of variation. It was a smooth random function of log
frequency (0.5--16\,kHz), band-limited to ripple densities below 0.5 ripples
per octave with an rms of 3\,dB (median peak-to-trough depth 10\,dB), applied as
a zero-phase gain. The cutoff follows Macpherson and Middlebrooks
% \citet{macpherson_vertical-plane_2003}
, who found that localization was disrupted by spectral ripple at densities of
0.5--2 ripples per octave but not by coarser variation, and concluded that
broad-scale spectral structure is discounted despite contributing most to the
shape of directional transfer functions. Each trial was
therefore strongly recolored while the 0.5--2 ripples/octave band, in which the
pinna notch shifts with elevation, was left largely untouched. The manipulation
thus tests whether what was learned during training is tied to the spectrum of
the training stimulus or generalizes across source spectra. It does not
separate a perceptual remapping from an explicitly learned association: a
listener who had learned to read the preserved band would succeed under either
account. Timing, bandwidth and level were unchanged, and the envelope of every
trial was stored.

Across 400 realizations the variable stimulus varied with an SD of 2.4\,dB,
against 0.4\,dB for the fixed-spectrum stimulus. Within the 0.5--2
ripples/octave band the elevation-dependent variation of the participants'
head-related transfer functions (SD across elevation at fixed azimuth) exceeded
the trial-to-trial variation of the source spectrum by a factor of 7.5; below
0.5 ripples/octave the two were comparable (factor 1.0).

% NB: every number in this file scales linearly with the envelope rms, so
% changing the depth is a proportional edit throughout. Values at the earlier
% 8 dB setting, for reference: 27 dB median peak-to-trough, 6.4 dB stimulus SD,
% 2.8 in the cue band, 0.4 below it. Settle this against the self-test
% (protocols/dev/stimulus_check.py cells 7-9) before submission.

\subsubsection*{Choice of envelope depth}

Envelope depth is bounded from both sides, and the two bounds are what fixed
it. It must be large enough that the coarse-scale coloration of the source is
no longer a reliable guide to elevation, and small enough that the elevation
cue itself survives. The upper bound is the more concrete of the two.
Macpherson and Middlebrooks
% \citet{macpherson_vertical-plane_2003}
report degradation of vertical-plane localization at 1 ripple per octave only
for peak-to-trough depths at or above 20\,dB, and at coarser densities even
40\,dB left performance largely intact; we therefore treated 20\,dB
peak-to-trough as a ceiling and stayed well below it. Two further
considerations argue for the same direction. The band limit is not sharp---a
function band-limited on a finite log-frequency axis has skirts---so a deeper
envelope leaks into the 0.5--2 ripples/octave region even though no energy is
assigned there, and the depth of the elevation cue itself varies between ears
and listeners, so a setting chosen against one head-related transfer function
must leave headroom for shallower ones.

The lower bound is set by the purpose of the manipulation. In the fixed-spectrum
condition the sound at the eardrum carries a near one-to-one map from absolute
spectrum to elevation, so a listener can in principle solve the task by
matching the overall spectrum to a stored template, without separating source
from filter. The envelope must make that route unreliable, which requires its
trial-to-trial variation to be at least comparable to the elevation-dependent
variation of the head-related transfer function at the same spectral scales. At
3\,dB rms the two are of the same magnitude below 0.5 ripples per octave
(factor 1.0), so overall coloration no longer specifies elevation, while the
cue band retains a 7.5:1 margin. Deeper envelopes would make the source
dominate rather than merely match at coarse scales, but only at the cost of
the margin in the cue band, and of exceeding the depth at which disruption has
been reported.

Variation was confined to the coarse band rather than spread across all
spectral scales for the same reason: variation inside the cue band would
contaminate the very structure being read, and contributes nothing to defeating
a template-matching strategy that the coarse band does not already achieve. It
was retained as an adjustable parameter and set to zero throughout, so that
deliberately contaminating the cue band remains available as a manipulation
check.

Finally, the parameters were verified behaviorally rather than only
acoustically. [PLACEHOLDER — self-test, cells 7-9: own unmodified HRTF,
binaural, fixed-spectrum vs variable-spectrum blocks; report elevation gain for
each and the difference. With intact ears and an unmodified transfer function,
a drop in elevation gain can only reflect the envelope encroaching on the cue,
so an absent or small drop is the direct evidence that the envelope is
transparent to localization.]

% --- supplement ------------------------------------------------------------
% The envelope is a discrete cosine series in dB on a 128-point log-frequency
% grid, coefficient k corresponding to k/10 ripples per octave; the constant
% term and all coefficients at or above 0.5 ripples/octave are zero, a zero
% coefficient contributing 0 dB, i.e. unity gain. Burst-to-burst correlation of
% the detrended log spectra within a trial was r = 0.95, versus r = -0.01 for
% the unmodified stimulus.

% --- what this manipulation can and cannot decide ---------------------------
% Rules out: a mapping tied to the overall spectrum of the training stimulus
%   ("this timbre = that elevation"), since gross coloration now varies trial to
%   trial by as much as elevation itself varies it. Note the strength of this
%   claim depends on the depth: at 3 dB the source only MATCHES the cue at coarse
%   scales, it does not swamp it, so the timbre route is made unreliable rather
%   than uninformative. Phrase it that way in the Discussion; "uninformative"
%   would need the deeper envelope, which M&M's depth threshold rules out.
% Does NOT rule out: a lookup restricted to the preserved 0.5-2 rip/oct band.
%   The pinna notch is a shape-within-band feature, so a remapping of that
%   structure onto space and a learned association between that structure and a
%   response both survive the manipulation. Arguably the two are not
%   distinguishable acoustically at all -- adaptation at the spectral-to-spatial
%   mapping stage IS the learning of new pattern-location associations.
% What would separate perceptual from explicit/procedural learning:
%   - response latency and dual-task cost (explicit strategies are slow and
%     interference-prone)
%   - generalization to untrained locations WITHIN the trained ear: a graded map
%     should interpolate, a lookup should not (cf. condition B)
%   - graded vs catastrophic degradation under band-limiting of the cue region
%   - aftereffects / coexistence on return to the native HRTF
%   - subjective report: does it sound like it comes from there, or is the
%     response reasoned to
% On whether the listener CAN discount the source spectrum here — two reasons
% they can, so a failure is informative rather than a ceiling of the paradigm:
%  (1) Presentation is not strictly monaural. The other ear receives the same
%      source convolved with a cepstrally smoothed (n_keep=4) DTF, so the
%      interaural log-difference cancels the source spectrum exactly while
%      retaining almost all of the trained ear's fine structure. The reference
%      channel is in fact smoother than in natural binaural listening, where
%      both ears carry notches. Its level falls with lateral angle (up to ~9 dB
%      down at 35 deg), so the reference degrades toward the periphery.
%  (2) Source-spectrum discounting is anyway achievable monaurally: models based
%      on the first/second spectral derivative recover direction for sources
%      with locally flat or locally constant-slope spectra (Zakarauskas &
%      Cynader 1993), and current sagittal-plane models use monaural spectral
%      gradients (Baumgartner et al. 2014, already cited in the Intro).
%
% Framing for the Discussion: Macpherson & Middlebrooks show the NORMAL system
% already discounts broad-scale spectral variation. The open question is whether
% a NEWLY acquired spectral-to-spatial mapping inherits that invariance, or is
% initially stimulus-specific and only becomes source-invariant with further
% exposure. That, not "sensory vs cognitive", is what the ripple blocks test.

% --- open items ------------------------------------------------------------
% - Citations to add once a .bib exists:
%     Macpherson & Middlebrooks (2003) JASA 114(1):430-445,
%       doi:10.1121/1.1582174  -- ripple-spectrum probe, the 0.5-2 rip/oct window
%     Zakarauskas & Cynader (1993) JASA 94(3):1323-1331 -- spectral-derivative
%       model, monaural source-spectrum invariance
%     Middlebrooks (1992), Baumgartner et al. (2014) -- already cited in Intro
% - DONE: depth is now reported peak-to-trough as well as rms, and the depth
%   bound is argued explicitly ("Choice of envelope depth").
% - Cue depth varies by ear and subject; the band ratios are scaled from FS's
%   modified left ear at az -20. Decide whether to report a group range instead
%   -- with a shallow envelope the cue-band margin is comfortable everywhere, so
%   a group minimum would be the more convincing number to give.
% - The self-test paragraph is a placeholder until cells 7-9 have been run.
%   If the fixed vs variable difference in elevation gain is not small, the
%   depth argument above has to be rewritten around whatever value survives.
% - FS's two "USO" blocks were run with the old generator (one fixed base
%   texture per session, 0.8 dB variation). They are NOT a variable-spectrum
%   test and should be reported as such, or dropped.
